\chapter{X.509, chuỗi tin cậy và cơ chế xác thực}
\label{chap:x509-xac-thuc}
\section{Cấu trúc chứng thư X.509 v3}
\label{sec:x509-structure}

Chuẩn ITU-T X.509 phiên bản 3 (RFC~5280~\cite{rfc5280}) định nghĩa định dạng chứng thư số bằng cú pháp ASN.1 (Abstract Syntax Notation One). Cấu trúc tổng thể của chứng thư X.509 v3 gồm ba trường chính ở cấp cao nhất:
\begin{verbatim}
Certificate  ::=  SEQUENCE  {
    tbsCertificate       TBSCertificate,
    signatureAlgorithm   AlgorithmIdentifier,
    signatureValue       BIT STRING  }
\end{verbatim}

Trong đó:
\begin{itemize}
    \item \texttt{tbsCertificate} (To-Be-Signed Certificate): Khối dữ liệu chứa toàn bộ thông tin định danh và tham số kỹ thuật cần được ký bảo đảm.
    \item \texttt{signatureAlgorithm}: Định danh thuật toán mật mã mà CA sử dụng để tạo chữ ký số (ví dụ: \texttt{sha256WithRSAEncryption} hoặc \texttt{ecdsa-with-SHA256}).
    \item \texttt{signatureValue}: Giá trị chữ ký số thực tế (chuỗi bit \texttt{BIT STRING}) do CA tính toán trên cấu trúc \texttt{tbsCertificate} mã hóa theo quy tắc chuẩn hóa DER. Việc sửa đổi dù chỉ một bit trên chứng thư đã phát hành sẽ phá hủy tính hợp lệ toán học của chữ ký và bị từ chối ngay lập tức (phép thử tính toàn vẹn này được kiểm chứng tại ca kiểm thử TC05).
\end{itemize}

Khối \texttt{tbsCertificate} là đối tượng cốt lõi phản ánh toàn bộ nội dung quản trị và kỹ thuật của chứng thư. Theo Mục~4.1 và 4.2 của RFC~5280~\cite{rfc5280}, các trường chuẩn và phần mở rộng thiết yếu bao gồm:

\subsection{Các trường dữ liệu chuẩn (Standard Fields)}
\label{subsec:standard-fields}

\begin{itemize}
    \item \textbf{Số sê-ri (Serial Number):} Số nguyên dương duy nhất do CA gán cho mỗi chứng thư ban hành, độ dài không quá 20 octet theo RFC~5280 (Mục 4.1.2.2). Để chống tấn công va chạm (collision) trên hàm băm, theo BR~\cite{cabforum_br}, CA công cộng phải sinh serial có ít nhất 64 bit ngẫu nhiên thay vì sử dụng cơ chế tăng tuần tự đơn giản.

    \item \textbf{Tên đơn vị cấp phát (Issuer):} Định danh CA đã ký và phát hành chứng thư dưới dạng Tên phân biệt (Distinguished Name -- DN) theo chuẩn X.500, gồm chuỗi các thuộc tính như Quốc gia ($C$), Tổ chức ($O$), Tên phổ biến ($CN$). Đối với chứng thư tự ký của Root CA, trường \texttt{issuer} trùng khớp hoàn toàn với trường \texttt{subject}.

    \item \textbf{Thời hạn hiệu lực (Validity):} Xác định khoảng thời gian mà CA cam kết duy trì thông tin trạng thái chứng thư, xác định bởi cặp mốc \texttt{notBefore} và \texttt{notAfter} (dùng \texttt{UTCTime} trước năm 2050 và \texttt{GeneralizedTime} từ năm 2050 trở đi). Thuật toán xác minh yêu cầu kiểm tra thời điểm xác minh nằm trong khoảng này; kết quả phụ thuộc chính sách của client (kiểm chứng tại ca kiểm thử TC03).

    \item \textbf{Tên chủ thể (Subject):} Định danh thực thể sở hữu khóa công khai bên trong chứng thư dưới dạng DN. RFC~5280 cho phép Subject rỗng khi danh tính nằm hoàn toàn trong SAN (khi đó SAN phải critical); chứng thư TLS hiện đại định danh tên máy chủ chủ yếu qua SAN.

    \item \textbf{Thông tin khóa công khai của chủ thể (Subject Public Key Info):} Chứa khóa công khai thực tế của chủ thể cùng mã OID định danh thuật toán và các tham số toán học liên quan (như modulus $N$, số mũ $e$ của RSA, hoặc định danh đường cong elip của ECDSA). Trường này cho phép relying party trích xuất đúng khóa và thuật toán để xác thực chữ ký của chủ thể mà không cần biết trước cấu hình.
\end{itemize}

\subsection{Các phần mở rộng X.509 v3 quan trọng (Standard Extensions)}
\label{subsec:standard-extensions}

Phần mở rộng v3 cho phép bổ sung các thuộc tính kiểm soát chính sách và bảo mật linh hoạt. Cấu trúc ASN.1 chuẩn của mỗi phần mở rộng gồm:
\begin{verbatim}
Extension ::= SEQUENCE {
    extnID      OBJECT IDENTIFIER,
    critical    BOOLEAN DEFAULT FALSE,
    extnValue   OCTET STRING }
\end{verbatim}

Mỗi extension gồm mã định danh (\texttt{extnID}), cờ \texttt{critical} và giá trị dữ liệu (\texttt{extnValue}). Nếu \texttt{critical: TRUE}, validator bắt buộc phải nhận diện và xử lý extension đó; nếu không hiểu, validator bắt buộc phải từ chối chứng thư. Các phần mở rộng cốt lõi bao gồm:
\begin{itemize}
    \item \textbf{Subject Alternative Name (SAN -- OID \texttt{2.5.29.17}):} Cho phép gán thêm nhiều danh tính cho chủ thể dưới các định dạng như tên miền (\texttt{dNSName}) hoặc địa chỉ IP (\texttt{iPAddress}). Các client hiện đại đối chiếu tên máy chủ với SAN, không dựa vào CN (kiểm chứng tại ca kiểm thử TC04).

    \item \textbf{Ràng buộc cơ bản -- Basic Constraints (OID \texttt{2.5.29.19}):} Xác định thực thể có quyền đóng vai trò CA để ký cấp chứng thư hay không qua cờ \texttt{cA} và giới hạn độ sâu tùy chọn \texttt{pathLenConstraint}. Chứng thư CA phải có extension này và phải critical (RFC~5280 mục 4.2.1.9). Thực thể cuối có thể không có extension này hoặc có \texttt{cA:FALSE}; extension tự nó không ngăn chặn, mà validator từ chối dùng chứng thư \texttt{cA:FALSE} làm CA (ràng buộc này ảnh hưởng trực tiếp đến việc xây dựng chuỗi ở ca TC06).

    \item \textbf{Mục đích sử dụng khóa -- Key Usage (OID \texttt{2.5.29.15}):} Chuỗi bit giới hạn các thao tác mật mã cơ sở của khóa công khai (như \texttt{digitalSignature}, \texttt{keyCertSign}, \texttt{cRLSign}), thường đặt \texttt{critical: TRUE} để ngăn ngừa dùng một khóa cho các ngữ cảnh mâu thuẫn. Cần lưu ý rằng \texttt{keyEncipherment} không cần cho TLS 1.3/ECDHE (chỉ cần \texttt{digitalSignature}).

    \item \textbf{Mục đích sử dụng khóa mở rộng (EKU -- OID \texttt{2.5.29.37}):} Xác định giao thức ứng dụng cấp cao mà khóa được áp dụng qua danh sách OID, tiêu biểu như \texttt{serverAuth} (xác thực máy chủ TLS) hoặc \texttt{codeSigning} (ký mã nguồn). Kết nối TLS máy chủ bị từ chối theo chính sách xác minh khi EKU có mặt mà không chứa \texttt{serverAuth} (kiểm chứng tại ca TC08). Chứng thư không có EKU thường không bị giới hạn mục đích.
\end{itemize}

\subsection{Trích xuất thực tế cấu trúc chứng thư trong môi trường Lab}
\label{subsec:lab-cert-extraction}

Dưới đây là kết quả giải mã cấu trúc chi tiết của chứng thư máy chủ TLS (\texttt{server-a.crt}) được cấp bởi Intermediate CA trong môi trường lab bằng lệnh:
\begin{verbatim}
$ openssl x509 -in demo/out/servers/server-a.crt -text -noout
\end{verbatim}

\begin{figure}[!htbp]
    \centering
    \includegraphics[width=0.85\textwidth, height=0.78\textheight, keepaspectratio]{openssl-x509-terminal-real.png}
    \caption{Kết quả thực thi lệnh trích xuất cấu trúc chứng thư X.509 v3 của máy chủ TLS (\texttt{server-a.crt}) trong môi trường thực nghiệm.}
    \label{fig:openssl-x509-terminal-real}
\end{figure}

\noindent\textbf{Phân tích từ kết quả thực nghiệm:}

Từ kết quả trích xuất trên Hình~\ref{fig:openssl-x509-terminal-real}, các trường cấu trúc của chứng thư máy chủ phản ánh đúng các thiết lập chính sách đã nghiên cứu:
\begin{itemize}
    \item \textbf{Phiên bản và Số sê-ri:} Chứng thư mang phiên bản \texttt{3 (0x2)} cho phép sử dụng các extension v3, và số sê-ri \texttt{4096 (0x1000)} được quản lý trong cơ sở dữ liệu \texttt{index.txt} của Intermediate CA.
    \item \textbf{Cặp định danh Issuer và Subject:} Thể hiện rõ mô hình phân tầng: \texttt{Issuer} trỏ về \texttt{CN = Lab Intermediate CA}, còn \texttt{Subject} định danh máy chủ dịch vụ \texttt{CN = server-a.test}.
    \item \textbf{SAN:} Chỉ định tên miền hợp lệ \texttt{DNS:server-a.test}, đáp ứng yêu cầu đối chiếu tên máy chủ của client (đối chứng với ca kiểm thử TC04).
    \item \textbf{Basic Constraints:} Thiết lập \texttt{critical, CA:FALSE}, định vị đây là chứng thư thực thể cuối (End-entity) và không có quyền ký phát chứng thư phái sinh.
    \item \textbf{Key Usage và EKU:} Bật \texttt{Digital Signature}, \texttt{Key Encipherment}; \texttt{EKU} kích hoạt \texttt{serverAuth}, đáp ứng điều kiện xác thực máy chủ TLS (đối chứng ca TC08).
    \item \textbf{Định danh khóa (SKI và AKI):} Trường SKI và AKI chứa mã băm nhận diện khóa công khai của máy chủ và của Intermediate CA, hỗ trợ thuật toán xây dựng và neo chuỗi tin cậy.
\end{itemize}

\section{Xây dựng và xác minh đường dẫn}
\label{sec:path-validation}

Một chứng thư máy chủ không tự mang đủ bằng chứng để được tin cậy. Relying party phải chỉ ra được một đường dẫn chứng thư (certification path) nối chứng thư đó tới một trust anchor mà chính nó đã cấu hình. Việc này gồm hai giai đoạn: xây dựng đường dẫn (path building) và xác minh đường dẫn (path validation, RFC~5280~\cite{rfc5280}, mục~6).

\subsection{Chuỗi chứng thư trong lab}
\label{subsec:lab-trust-chain}

Lab dùng chuỗi ba mắt xích $\mathrm{Root\ CA} \rightarrow \mathrm{Intermediate\ CA} \rightarrow \text{chứng thư máy chủ}$, như minh họa ở Hình~\ref{fig:trust-chain-validation}.
\begin{itemize}
    \item \textbf{Root CA} (\texttt{Lab Root CA}): Chứng thư tự ký ($\mathrm{Issuer} \equiv \mathrm{Subject}$), \texttt{basicConstraints} critical với \texttt{CA:TRUE}. Khóa $SK_{\mathrm{Root}}$ dùng để ký chứng thư của Intermediate CA.
    \item \textbf{Intermediate CA} (\texttt{Lab Intermediate CA}): Được Root ký phát với chữ ký $\sigma_{\mathrm{Sub}} = \mathrm{Sign}_{SK_{\mathrm{Root}}}(\mathrm{TBS}_{\mathrm{Sub}})$. Chứng thư này có cờ \texttt{CA:TRUE} và \texttt{pathlen:0}, nghĩa là chỉ được cấp chứng thư thực thể cuối, không được cấp thêm CA cấp dưới. Khóa $SK_{\mathrm{Sub}}$ dùng để ký chứng thư máy chủ.
    \item \textbf{Chứng thư máy chủ} (\texttt{server-a.test}): Được Intermediate CA ký phát với chữ ký $\sigma_{\mathrm{Server}} = \mathrm{Sign}_{SK_{\mathrm{Sub}}}(\mathrm{TBS}_{\mathrm{Server}})$, thiết lập \texttt{CA:FALSE}, SAN \texttt{DNS:server-a.test} và EKU \texttt{serverAuth}.
\end{itemize}

\clearpage

\begin{figure}[!htbp]
\centering
\resizebox{0.88\textwidth}{!}{%
\begin{tikzpicture}[
    % Định dạng các khối theo sắc xám học thuật nhất quán với Hình 3.1
    rootbox/.style={draw=black, very thick, double, fill=gray!18, rounded corners=6pt, text width=7.2cm, align=left, inner sep=5pt, font=\small, text=black},
    cabox/.style={draw=black, thick, fill=gray!10, rounded corners=6pt, text width=7.2cm, align=left, inner sep=5pt, font=\small, text=black},
    leafbox/.style={draw=black, thick, fill=white, rounded corners=6pt, text width=7.2cm, align=left, inner sep=5pt, font=\small, text=black},
    storebox/.style={draw=black, thick, dashed, fill=gray!6, rounded corners=6pt, text width=6.6cm, align=center, inner sep=5pt, font=\small, text=black},
    % Kiểu mũi tên nét liền (tiến trình ký) và nét đứt (tiến trình thẩm định)
    issuancearrow/.style={-{Stealth[scale=1.05]}, thick, black},
    valarrow/.style={-{Stealth[scale=1.1]}, very thick, dash pattern=on 5pt off 2.5pt, black},
    trustarrow/.style={-{Stealth[scale=1.1]}, very thick, dash pattern=on 6pt off 3pt, black}
]

    % --- KHỐI GỐC TIN CẬY (TRUST ANCHOR TRÊN CÙNG) ---
    \node[storebox] (store) at (0, 15.2) {
        \textbf{Kho chứng thư tin cậy (Root Trust Store)}\\[1pt]
        {\scriptsize Quản lý bởi Hệ điều hành / Trình duyệt (Relying Party)\\[1pt]Chứa bản sao khóa công khai $PK_{\mathrm{Root}}$ được tin cậy mặc định}
    };

    % --- KHỐI 1: Root CA ---
    \node[rootbox] (root) at (0, 10.4) {
        \textbf{Mắt xích 1: Root CA (Trust Anchor)}\\[1pt]
        {\scriptsize \texttt{Subject: CN = Lab Root CA}}\\[1pt]
        {\scriptsize \texttt{Issuer:  CN = Lab Root CA} (Tự ký - Self-signed)}\\[1pt]
        {\scriptsize Khóa: Cặp khóa gốc $(SK_{\mathrm{Root}}, PK_{\mathrm{Root}})$}\\[1pt]
        {\scriptsize \texttt{Basic Constraints: critical, CA:TRUE}}\\[1pt]
        {\scriptsize \texttt{SKI:} Mã băm nhận diện khóa $PK_{\mathrm{Root}}$}
    };

    % --- KHỐI 2: Intermediate CA ---
    \node[cabox] (subca) at (0, 5.2) {
        \textbf{Mắt xích 2: Intermediate CA (Sub-CA)}\\[1pt]
        {\scriptsize \texttt{Subject: CN = Lab Intermediate CA}}\\[1pt]
        {\scriptsize \texttt{Issuer:  CN = Lab Root CA}}\\[1pt]
        {\scriptsize Khóa: Cặp khóa trung gian $(SK_{\mathrm{Sub}}, PK_{\mathrm{Sub}})$}\\[1pt]
        {\scriptsize \texttt{Basic Constraints: critical, CA:TRUE, pathlen:0}}\\[1pt]
        {\scriptsize \texttt{AKI:} Trỏ tới $\mathrm{SKI}_{\mathrm{Root}}$ \textbar{} \texttt{SKI:} Băm khóa $PK_{\mathrm{Sub}}$}
    };

    % --- KHỐI 3: Server Certificate ---
    \node[leafbox] (server) at (0, 0.0) {
        \textbf{Mắt xích 3: Server Certificate (Thực thể cuối)}\\[1pt]
        {\scriptsize \texttt{Subject: CN = server-a.test}}\\[1pt]
        {\scriptsize \texttt{Issuer:  CN = Lab Intermediate CA}}\\[1pt]
        {\scriptsize Khóa: Khóa công khai $PK_{\mathrm{Server}}$ ($SK_{\mathrm{Server}}$ lưu an toàn trên Server)}\\[1pt]
        {\scriptsize \texttt{Basic Constraints: CA:FALSE} \textbar{} \texttt{EKU: serverAuth}}\\[1pt]
        {\scriptsize \texttt{SAN: DNS:server-a.test} \textbar{} \texttt{AKI:} Trỏ tới $\mathrm{SKI}_{\mathrm{Sub}}$}
    };

    % Neo Root CA vào Trust Store
    \draw[trustarrow] (root.north) -- (store.south)
        node[midway, right=8pt, align=left, font=\footnotesize\bfseries] {Neo niềm tin\\[1pt](Trust Anchor)};

    % ==============================================================
    % TIÊU ĐỀ HAI TIẾN TRÌNH Ở ĐỈNH (KHÔNG XUNG ĐỘT)
    % ==============================================================
    \node[align=center, font=\footnotesize\bfseries, text=black] at (-6.3, 15.2) {
        [TIẾN TRÌNH KÝ PHÁT]\\[2pt]
        {\scriptsize (Dùng khóa bí mật $SK$)}
    };

    \node[align=center, font=\footnotesize\bfseries, text=black] at (6.5, 15.2) {
        [TIẾN TRÌNH THẨM ĐỊNH]\\[2pt]
        {\scriptsize (Dùng khóa công khai $PK$)}
    };

    % ==============================================================
    % TIẾN TRÌNH KÝ PHÁT (BÊN TRÁI - TỪ TRÊN XUỐNG)
    % ==============================================================
    % Root CA ký Sub-CA
    \draw[issuancearrow] (root.west) -- (-4.5, 10.4) -- (-4.5, 5.2) -- (subca.west);
    \node[anchor=east, align=right, font=\footnotesize] at (-4.7, 7.8) {
        \textbf{1. Ký phát Sub-CA:}\\[1pt]
        $\sigma_{\mathrm{Sub}} = \mathrm{Sign}_{SK_{\mathrm{Root}}}(\mathrm{TBS}_{\mathrm{Sub}})$\\[1pt]
        {\scriptsize $SK_{\mathrm{Root}}$ ký lên chứng thư Sub-CA}
    };

    % Sub-CA ký Server Certificate
    \draw[issuancearrow] (subca.west) -- (-4.5, 5.2) -- (-4.5, 0.0) -- (server.west);
    \node[anchor=east, align=right, font=\footnotesize] at (-4.7, 2.6) {
        \textbf{2. Ký phát Server:}\\[1pt]
        $\sigma_{\mathrm{Server}} = \mathrm{Sign}_{SK_{\mathrm{Sub}}}(\mathrm{TBS}_{\mathrm{Server}})$\\[1pt]
        {\scriptsize $SK_{\mathrm{Sub}}$ ký lên chứng thư Server}
    };

    % ==============================================================
    % TIẾN TRÌNH THẨM ĐỊNH RFC 5280 (BÊN PHẢI - TỪ DƯỚI LÊN)
    % ==============================================================
    % Thẩm định Server Certificate
    \draw[valarrow] (server.east) -- (4.7, 0.0) -- (4.7, 5.2) -- (subca.east);
    \node[anchor=west, align=left, font=\footnotesize] at (4.9, 2.6) {
        \textbf{3. Thẩm định Server Certificate:}\\[1pt]
        $\mathrm{Verify}_{PK_{\mathrm{Sub}}}(\sigma_{\mathrm{Server}}) = \mathrm{True}$\\[1pt]
        {\scriptsize $\bullet$ Đối chiếu: $\mathrm{Issuer} \equiv \mathrm{Subject}_{\mathrm{Sub}}$}\\[1pt]
        {\scriptsize $\bullet$ Khớp khóa: $\mathrm{AKI}_{\mathrm{Server}} \equiv \mathrm{SKI}_{\mathrm{Sub}}$}
    };

    % Thẩm định Intermediate CA
    \draw[valarrow] (subca.east) -- (4.7, 5.2) -- (4.7, 10.4) -- (root.east);
    \node[anchor=west, align=left, font=\footnotesize] at (4.9, 7.8) {
        \textbf{4. Thẩm định Intermediate CA:}\\[1pt]
        $\mathrm{Verify}_{PK_{\mathrm{Root}}}(\sigma_{\mathrm{Sub}}) = \mathrm{True}$\\[1pt]
        {\scriptsize $\bullet$ Đối chiếu: $\mathrm{Issuer} \equiv \mathrm{Subject}_{\mathrm{Root}}$}\\[1pt]
        {\scriptsize $\bullet$ Khớp khóa: $\mathrm{AKI}_{\mathrm{Sub}} \equiv \mathrm{SKI}_{\mathrm{Root}}$}
    };

\end{tikzpicture}%
}

\caption{Chuỗi chứng thư Root CA $\rightarrow$ Intermediate CA $\rightarrow$ server trong lab: quan hệ ký phát và quan hệ thẩm định giữa các mắt xích.}
\label{fig:trust-chain-validation}
\end{figure}

\subsection{Quy trình xây dựng và xác minh}
\label{subsec:path-building-validation}

\noindent\textbf{Xây dựng đường dẫn.} Xuất phát từ chứng thư máy chủ, relying party tìm chứng thư của bên phát hành bằng cách khớp trường \texttt{Issuer} của chứng thư dưới với \texttt{Subject} của chứng thư trên, đồng thời khớp trường AKI với SKI (chỉ khớp tên là chưa đủ, vì một CA có thể có nhiều khóa cùng tên). Các chứng thư trung gian thường lấy từ chuỗi do máy chủ gửi kèm và chưa được tin cậy ở thời điểm này. Quá trình dừng khi gặp một chứng thư nằm trong kho trust anchor của client.

\noindent\textbf{Trust anchor.} Một Root CA tự ký chỉ trở thành trust anchor khi client được cấu hình để tin cậy nó (trong lab là tùy chọn \texttt{-CAfile}). Độ tin cậy của Root đến từ cấu hình đó, không đến từ chữ ký tự ký, vì bất kỳ ai cũng tạo được chứng thư tự ký cho chính mình.

\noindent\textbf{Xác minh đường dẫn.} Thuật toán của RFC~5280 quy định quy trình xử lý từng chứng thư trên đường dẫn từ trust anchor xuống chứng thư thực thể cuối qua 6 bước kiểm tra tuần tự:
\begin{enumerate}\setlength{\itemsep}{1pt}\setlength{\parskip}{0pt}
    \item \textbf{Xây dựng đường dẫn (Path Construction):} Truy vết các mắt xích từ chứng thư máy chủ ngược lên Trust Anchor dựa trên cặp định danh Issuer/Subject và AKI/SKI.
    \item \textbf{Xác minh chữ ký số (Signature Verification):} Dùng khóa công khai của bên cấp phát để kiểm tra tính toàn vẹn toán học của chữ ký trên từng chứng thư trong chuỗi.
    \item \textbf{Kiểm tra khung thời hạn hiệu lực (Validity Check):} Đối chiếu thời điểm xác minh với khoảng \texttt{[notBefore, notAfter]} của từng mắt xích.
    \item \textbf{Kiểm tra ràng buộc CA và extension liên quan (Constraints Check):} Thẩm tra \texttt{basicConstraints} (cờ \texttt{cA:TRUE} cho CA, giới hạn \texttt{pathLen}) và các ràng buộc chính sách.
    \item \textbf{Kiểm tra danh tính máy chủ và mục đích sử dụng (Identity \& Purpose Matching):} So khớp tên miền truy cập với SAN và đối chiếu cờ EKU (\texttt{serverAuth}) theo yêu cầu ứng dụng.
    \item \textbf{Kiểm tra trạng thái thu hồi (Revocation Status):} Tra cứu số sê-ri trong CRL hoặc truy vấn OCSP responder nếu chính sách kiểm tra thu hồi được kích hoạt.
\end{enumerate}

Cần lưu ý quan trọng rằng công cụ \texttt{openssl verify} không mặc định thực thi toàn bộ các bước trên. Hành vi xác minh phụ thuộc trực tiếp vào các tùy chọn dòng lệnh: kiểm tra hostname bắt buộc phải khai báo tường minh qua \texttt{-verify\_hostname}; kiểm tra mục đích TLS server cần chỉ định \texttt{-purpose sslserver}; và kiểm tra thu hồi đòi hỏi cờ \texttt{-crl\_check}. Nếu thiếu các tùy chọn này, lệnh chỉ thực hiện xác minh chuỗi cơ bản. Mỗi mắt xích phải đạt các điều kiện ở Bảng~\ref{tab:path-validation-conditions}. Chỉ cần một điều kiện không đạt là chứng thư bị từ chối theo chính sách xác minh được bật, dù các chữ ký vẫn đúng về mặt toán học. Đây là cơ sở của giả thuyết $H_1$ và được kiểm chứng bằng các ca TC02--TC08, TC10 ở Chương~\ref{chap:experiment}.

\begin{table}[!htbp]
\centering
\small
\renewcommand{\arraystretch}{1.25}
\begin{tabularx}{\textwidth}{>{\raggedright\arraybackslash}X>{\raggedright\arraybackslash}p{4.2cm}>{\raggedright\arraybackslash}p{2.8cm}}
\toprule
\textbf{Điều kiện kiểm tra} & \textbf{Áp dụng cho} & \textbf{Ca kiểm thử} \\
\midrule
Chữ ký của chứng thư được tạo bằng khóa của bên phát hành & Mọi chứng thư, trừ trust anchor & TC05 \\
\midrule
Chuỗi nối được tới trust anchor đã cấu hình & Toàn đường dẫn & TC02, TC06, TC10 \\
\midrule
Thời điểm xác minh nằm trong [\texttt{notBefore}, \texttt{notAfter}] & Mọi chứng thư & TC03 \\
\midrule
Ràng buộc CA: \texttt{CA:TRUE} với mắt xích trung gian, \texttt{pathLenConstraint} không bị vượt & Root, Intermediate & Không có ca riêng \\
\midrule
Mục đích sử dụng: Key Usage / EKU phù hợp (\texttt{serverAuth}) & Chứng thư máy chủ & TC08 \\
\midrule
Tên máy chủ khớp với SAN & Chứng thư máy chủ & TC04 \\
\midrule
Trạng thái thu hồi (chỉ khi được bật) & Mọi chứng thư & TC07 (xem Mục~\ref{sec:revocation-mechanisms}) \\
\bottomrule
\end{tabularx}
\caption{Các điều kiện xác minh đường dẫn và ca kiểm thử tương ứng}
\label{tab:path-validation-conditions}
\end{table}


Khi mọi điều kiện đều đạt, kết quả là PASS (TC01, TC09).

\newpage

\subsection{Quan sát từ lab}
\label{subsec:lab-observations}

Hình~\ref{fig:openssl-verify-terminal} trình bày hai lần chạy \texttt{openssl verify} với cùng chứng thư \texttt{server-a.crt}.

\begin{figure}[!htbp]
\centering
\includegraphics[width=0.90\textwidth]{openssl-verify-terminal.png}
\caption{Kết quả \texttt{openssl verify} với chuỗi đầy đủ (OK) và khi thiếu chứng thư Intermediate CA (lỗi 20, TC06).}
\label{fig:openssl-verify-terminal}
\end{figure}

\noindent\textbf{Lần 1 (đủ chuỗi):} Khi cung cấp Root qua \texttt{-CAfile} và Intermediate qua \texttt{-untrusted}, kết quả là \texttt{OK}, kèm \texttt{-purpose sslserver} và \texttt{-verify\_hostname server-a.test} nên EKU và SAN cũng được kiểm. Chuỗi in ra gồm ba mức: depth~0 là server, depth~1 là Lab Intermediate CA, depth~2 là Lab Root CA. Nhãn \texttt{(untrusted)} ở depth~0 và 1 nghĩa là hai chứng thư này không nằm trong kho tin cậy; chỉ Root là trust anchor.

\noindent\textbf{Lần 2 (thiếu Intermediate):} Khi bỏ \texttt{-untrusted}, OpenSSL báo \texttt{error 20 at 0 depth lookup: unable to get local issuer certificate}. Chữ ký trên chứng thư máy chủ vẫn đúng, nhưng không tìm được chứng thư của bên phát hành nên không dựng được đường dẫn tới trust anchor. Lỗi gắn với depth~0 vì đó là chứng thư không tìm được issuer. Đây là ca TC06.

Cần lưu ý hai điểm. Thứ nhất, cả hai lệnh không bật kiểm tra thu hồi, nên \texttt{OK} chỉ cho biết các điều kiện trong Bảng~\ref{tab:path-validation-conditions} (trừ thu hồi) đã đạt; trạng thái thu hồi được xét riêng ở Mục~\ref{sec:revocation-mechanisms}. Thứ hai, kết quả của TC06 phụ thuộc vào client: OpenSSL không tự bổ sung chứng thư trung gian còn thiếu, trong khi một số client khác có cơ chế bù (qua AIA hoặc bộ nhớ đệm). Điều này phù hợp với giới hạn của lab đã nêu ở Mục~6.4.

\section{Thu hồi và trạng thái chứng thư}
\label{sec:revocation-mechanisms}

Một chứng thư có chữ ký đúng và còn trong khoảng \texttt{[notBefore, notAfter]} vẫn có thể không còn đáng tin, ví dụ khi khóa bí mật bị lộ hoặc CA cấp sai (Mục~\ref{sec:certificate-lifecycle}, Bước~9). Relying party vì vậy cần một cách hỏi chứng thư đã bị thu hồi chưa. RFC~5280~\cite{rfc5280} định nghĩa danh sách thu hồi chứng thư (CRL), còn RFC~6960~\cite{rfc6960} định nghĩa giao thức trạng thái chứng thư trực tuyến (OCSP). Hai cơ chế trả lời cùng một câu hỏi nhưng khác nhau ở cách thông tin được phát hành và lấy về.

\subsection{Danh sách thu hồi chứng thư (CRL)}
\label{subsec:crl}

CRL là danh sách có chữ ký do CA (hoặc CRL issuer được ủy quyền) phát hành định kỳ (RFC~5280~\cite{rfc5280}, mục~5). Mỗi CRL chứa số sê-ri các chứng thư bị thu hồi, thời điểm thu hồi và mã lý do tùy chọn (ví dụ \texttt{keyCompromise}), cùng hai trường \texttt{thisUpdate} và \texttt{nextUpdate} cho biết CRL được tạo khi nào và dự kiến được thay khi nào. Chứng thư chỉ nơi công bố CRL qua extension \texttt{CRL Distribution Points}. Khi xác minh, relying party tải CRL, kiểm tra chữ ký bằng khóa của bên phát hành, kiểm tra CRL còn hiệu lực rồi tra số sê-ri của chứng thư (RFC~5280~\cite{rfc5280}, mục~6.3).

Điểm mạnh là client có thể dùng bản CRL đã lưu đệm đến hết \texttt{nextUpdate}, nên không cần liên lạc với CA ở mỗi kết nối. Điểm yếu là độ trễ: thu hồi chỉ có hiệu lực với client khi CRL mới được phát hành và tải về. Kích thước CRL cũng tăng theo số chứng thư bị thu hồi (RFC~5280 có cơ chế delta CRL để giảm bớt).

\subsection{Giao thức trạng thái chứng thư trực tuyến (OCSP)}
\label{subsec:ocsp}

Với OCSP, client gửi tới OCSP responder một truy vấn chứa định danh của chứng thư (băm tên và khóa của bên phát hành cùng số sê-ri). Địa chỉ responder nằm trong extension \texttt{Authority Information Access} (RFC~6960~\cite{rfc6960}). Responder trả về phản hồi có chữ ký với một trong ba trạng thái \texttt{good}, \texttt{revoked} hoặc \texttt{unknown}, kèm khoảng thời gian phản hồi có hiệu lực. Phản hồi do CA hoặc một responder được CA ủy quyền ký, và truy vấn có thể kèm nonce để chống phát lại. Máy chủ cũng có thể đính kèm sẵn phản hồi OCSP trong bắt tay TLS (OCSP stapling) để client không phải tự truy vấn (xem Mục~\ref{sec:trust-models-comparison}).

Điểm mạnh là thông tin theo từng chứng thư, nhỏ gọn và mới hơn CRL. Điểm yếu là cần responder trực tuyến (khả dụng), responder biết client đang kiểm tra chứng thư nào (riêng tư), và khi không liên lạc được responder thì client phải tự quyết định chấp nhận hay từ chối.

Trong Web PKI, BR từ phiên bản 2.0.1 (hiệu lực 15/03/2024) đã chuyển OCSP thành tùy chọn và đặt yêu cầu công bố CRL~\cite{cabforum_br}. Phân tích sâu hơn ở Mục~\ref{sec:trust-models-comparison}.

\subsection{So sánh CRL và OCSP}
\label{subsec:crl-vs-ocsp}

Bảng~\ref{tab:crl-vs-ocsp} tổng hợp so sánh giữa hai cơ chế CRL và OCSP.

\begin{table}[!htbp]
\centering
\small
\renewcommand{\arraystretch}{1.2}
\begin{tabularx}{\textwidth}{>{\raggedright\arraybackslash}p{3.2cm}>{\raggedright\arraybackslash}X>{\raggedright\arraybackslash}X}
\toprule
\textbf{Tiêu chí} & \textbf{CRL} & \textbf{OCSP} \\
\midrule
\textbf{Đơn vị thông tin} & Danh sách số sê-ri bị thu hồi của một CA & Trạng thái của một chứng thư \\
\midrule
\textbf{Độ mới} & Giới hạn bởi chu kỳ phát hành (\texttt{nextUpdate}) & Mới hơn, tùy cách responder cập nhật \\
\midrule
\textbf{Băng thông} & Lớn khi danh sách dài & Nhỏ mỗi truy vấn nhưng nhiều truy vấn \\
\midrule
\textbf{Riêng tư} & Không lộ chứng thư cụ thể đang kiểm tra & Responder biết chứng thư đang được kiểm tra \\
\midrule
\textbf{Khả dụng} & Dùng được bản đã lưu đệm đến \texttt{nextUpdate} & Phụ thuộc responder trực tuyến (trừ khi có stapling) \\
\midrule
\textbf{Trong lab} & Triển khai và kiểm thử (TC07) & Chỉ phân tích lý thuyết \\
\bottomrule
\end{tabularx}
\caption{So sánh CRL và OCSP}
\label{tab:crl-vs-ocsp}
\end{table}


\subsection{Trạng thái \texttt{good} không đồng nghĩa với chứng thư hợp lệ}
\label{subsec:good-vs-valid}

RFC~6960 nêu rõ: \texttt{good} chỉ cho biết tối thiểu là chứng thư chưa bị thu hồi, không nhất thiết có nghĩa là chứng thư từng được cấp hay thời điểm phản hồi nằm trong thời hạn của chứng thư (RFC~6960~\cite{rfc6960}, mục~2.2). OCSP chỉ trả lời một câu hỏi trong Bảng~\ref{tab:path-validation-conditions} (trạng thái thu hồi) và không thay thế các câu còn lại: chuỗi tới trust anchor, \texttt{[notBefore, notAfter]}, Basic Constraints, EKU và SAN. Một chứng thư hết hạn, sai tên máy chủ hoặc sai EKU vẫn có thể nhận \texttt{good}. Nguyên lý này áp dụng cả cho CRL: không thấy số sê-ri trong CRL chỉ có nghĩa là chưa bị liệt kê.

Sự cố DigiNotar minh họa rủi ro này. Theo Fox-IT~\cite{prins2011diginotar}, hành vi thông thường theo RFC~2560 (tiền thân của RFC~6960) là trả \texttt{good} cho số sê-ri không có bản ghi. Vì vậy các chứng thư giả không có trong hồ sơ vẫn có thể qua bước kiểm tra thu hồi. Ngày 1/9/2011, responder được chuyển sang danh sách trắng, trả \texttt{revoked} cho mọi số sê-ri không xác định.

\noindent\textbf{Liên hệ lab.} Lab chỉ kiểm tra CRL: Intermediate CA thu hồi chứng thư bằng \texttt{openssl ca -revoke}, phát hành CRL bằng \texttt{openssl ca -gencrl}, và client xác minh với CRL được chỉ định tường minh. TC07 kiểm tra rằng chứng thư đã bị liệt kê bị từ chối khi kiểm tra CRL được bật; kết quả cụ thể ở Mục~\ref{sec:exp-results}. Lệnh \texttt{openssl verify} mặc định không kiểm tra thu hồi, nên với client không bật CRL, chứng thư bị thu hồi vẫn có thể qua các bước còn lại. Điều này khớp với giả thuyết $H_2$.

\section{Liên hệ TLS}
\label{sec:tls-relation}

Các mục trước xét việc xác minh chứng thư một cách độc lập. Trong thực tế, chứng thư được trình bày trong bắt tay TLS 1.3~\cite{rfc8446}, và việc một chứng thư ``được chấp nhận'' là kết quả của nhiều kiểm tra do các bên khác nhau thực hiện. Mục này làm rõ vai trò của chứng thư và thông điệp \texttt{CertificateVerify} trong bắt tay, rồi phân định kiểm tra nào thuộc giao thức, thư viện, ứng dụng hay chính sách riêng của trình duyệt.

\subsection{Chứng thư và \texttt{CertificateVerify} trong bắt tay TLS 1.3}
\label{subsec:tls-cert-verify}

Sau khi hai bên thỏa thuận khóa phiên (ClientHello và ServerHello mang \texttt{key\_share}), máy chủ gửi ba thông điệp xác thực (RFC~8446~\cite{rfc8446}, mục~4.4):
\begin{itemize}
    \item \textbf{\texttt{Certificate}:} Chuỗi chứng thư, chứng thư thực thể cuối đứng đầu, tiếp theo là các chứng thư trung gian. Chứng thư của trust anchor có thể được lược bỏ khi client đã có sẵn (RFC~8446~\cite{rfc8446}, mục~4.4.2). Đây là lý do thiếu Intermediate (TC06) khiến client không dựng được đường dẫn.
    \item \textbf{\texttt{CertificateVerify}:} Chữ ký của máy chủ, tạo bằng khóa bí mật tương ứng với khóa công khai trong chứng thư thực thể cuối, tính trên giá trị băm của toàn bộ các thông điệp bắt tay trước đó (RFC~8446~\cite{rfc8446}, mục~4.4.3).
    \item \textbf{\texttt{Finished}:} Mã xác thực (MAC) trên toàn bộ bắt tay, xác nhận hai bên cùng suy ra một bộ khóa và một bản ghi bắt tay (RFC~8446~\cite{rfc8446}, mục~4.4.4).
\end{itemize}

Hai chữ ký trong bắt tay có vai trò khác nhau và bổ sung cho nhau:
\begin{itemize}
    \item Chữ ký của CA trên chứng thư trả lời câu hỏi ``khóa này thuộc về tên này theo CA mà client tin cậy''.
    \item \texttt{CertificateVerify} trả lời câu hỏi ``bên đang kết nối thật sự nắm khóa bí mật tương ứng, trong chính phiên này'', vì nội dung được ký bao gồm cả \texttt{key\_share} của phiên.
\end{itemize}

Thiếu \texttt{CertificateVerify}, kẻ tấn công có thể trình lại chứng thư công khai của máy chủ mà không có khóa bí mật. Thiếu xác minh đường dẫn, kẻ tấn công trình chứng thư tự tạo kèm chữ ký đúng, đúng như kịch bản thay thế khóa của Eve ở Mục~\ref{subsec:eve-mitm-attack}. Cần cả hai thì kết nối mới chống được kịch bản đó.

Hai chữ ký này cũng khác nhau về thuật toán. Chữ ký \texttt{CertificateVerify} phải dùng các thuật toán client chấp nhận trong \texttt{signature\_algorithms}, và với RSA thì dùng RSA-PSS (RFC~8446~\cite{rfc8446}, mục~4.2.3). Chữ ký của CA trên chứng thư (ví dụ \texttt{sha256WithRSAEncryption} ở Hình~\ref{fig:openssl-x509-terminal-real}) có thể dùng PKCS\#1 v1.5.

Giới hạn quan trọng là \texttt{CertificateVerify} chỉ chứng minh sở hữu khóa, không chứng minh chứng thư được cấp đúng. Với DigiNotar, kẻ tấn công vừa có chứng thư giả \texttt{*.google.com} do CA bị chiếm ký, vừa nắm khóa bí mật tương ứng, nên mọi kiểm tra mật mã trong bắt tay đều đạt. Theo Fox-IT~\cite{prins2011diginotar}, sự cố lộ ra vì Chrome có kiểm tra bổ sung (certificate pinning) cho chứng thư của Google. Đây là cơ sở cho Mục~\ref{sec:valid-signature-limits}.

\subsection{Phân định các kiểm tra}
\label{subsec:tls-checks}

RFC~8446 chỉ quy định phần xác thực nằm trong giao thức. Việc xác minh chứng thư chi tiết được dẫn tới RFC~5280 (RFC~8446~\cite{rfc8446}, mục~4.4.2.4; \cite{rfc5280}). Việc đối chiếu tên máy chủ do RFC~9525~\cite{rfc9525} quy định, và RFC này chỉ xét tên trong chứng thư máy chủ cuối, không xét việc xác minh chuỗi. Bảng~\ref{tab:tls-checks} phân định ai chịu trách nhiệm kiểm tra gì.

\begin{table}[!htbp]
\centering
\small
\renewcommand{\arraystretch}{1.12}
\begin{tabularx}{\textwidth}{>{\raggedright\arraybackslash}p{3.5cm}>{\raggedright\arraybackslash}X>{\raggedright\arraybackslash}p{2.5cm}>{\raggedright\arraybackslash}p{2.3cm}}
\toprule
\textbf{Kiểm tra} & \textbf{Bên thực hiện} & \textbf{Nguồn quy định} & \textbf{Trong lab} \\
\midrule
Cấu trúc thông điệp \texttt{Certificate}, kiểu chứng thư, khóa tương thích & Giao thức TLS & RFC~8446~\cite{rfc8446} & Không kiểm tra \\
\midrule
\texttt{CertificateVerify} (sở hữu khóa) và \texttt{Finished} & Giao thức TLS & RFC~8446~\cite{rfc8446} & Không kiểm tra \\
\midrule
Dựng và xác minh đường dẫn, thời hạn, Basic Constraints, Key Usage/EKU & Thư viện TLS, với trust anchor, thời điểm và mục đích sử dụng do ứng dụng cấu hình & RFC~5280~\cite{rfc5280} & TC01--TC03, TC05, TC06, TC08, TC10 \\
\midrule
Trạng thái thu hồi & Ứng dụng quyết định có kiểm tra không, dùng CRL hay OCSP, và hành vi khi lỗi; TLS chỉ cho phép đính kèm phản hồi OCSP (stapling) & RFC~5280~\cite{rfc5280}, RFC~6960~\cite{rfc6960} & TC07 (chỉ CRL) \\
\midrule
Khớp tên máy chủ với SAN & Ứng dụng, hoặc thư viện khi ứng dụng yêu cầu & RFC~9525~\cite{rfc9525} & TC04 \\
\midrule
Chính sách riêng của trình duyệt và Root Store: danh mục trust anchor, yêu cầu CT, pinning, giới hạn thời hạn chứng thư & Trình duyệt hoặc nhà cung cấp Root Store & Ngoài RFC~5280 và 8446; một phần theo BR~\cite{cabforum_br} & Không có \\
\bottomrule
\end{tabularx}
\caption{Phân định các kiểm tra trong kết nối TLS 1.3}
\label{tab:tls-checks}
\end{table}


Hai hệ quả cần nhấn mạnh:
\begin{itemize}
    \item Nhiều thư viện không tự bật kiểm tra tên máy chủ hay kiểm tra thu hồi; ứng dụng phải yêu cầu (trong lab là \texttt{-verify\_hostname} và kiểm tra CRL được bật tường minh). Vì vậy một kết nối TLS ``thành công'' chưa chắc đã qua đủ các kiểm tra ở Bảng~\ref{tab:path-validation-conditions}.
    \item Cùng một chứng thư có thể được chấp nhận ở client này và bị từ chối ở client khác tùy cấu hình và chính sách. Các ca kiểm thử trong lab phản ánh tầng ``thư viện OpenSSL kèm cấu hình'' chứ không phải chính sách của trình duyệt (xem Mục~\ref{sec:exp-limitations}).
\end{itemize}

\noindent\textbf{Phạm vi của công cụ trong lab.} Lệnh \texttt{openssl verify} đọc các chứng thư được cung cấp và chạy thuật toán xác minh đường dẫn theo hồ sơ PKIX (RFC~5280~\cite{rfc5280}) tại một thời điểm xác định. Lệnh này không thiết lập bắt tay TLS, không trao đổi \texttt{key\_share}, không tính giá trị băm của bản ghi bắt tay và không kiểm tra việc nắm giữ khóa bí mật thông qua \texttt{CertificateVerify} (RFC~8446~\cite{rfc8446}). Vì vậy kết quả TC01--TC10 phản ánh tầng xác minh chứng thư, không phản ánh toàn bộ một kết nối TLS. Việc kiểm tra bằng \texttt{openssl s\_server} và \texttt{openssl s\_client} là hướng mở rộng.
